\IfFileExists{stacks-project.cls}{%
\documentclass{stacks-project}
}{%
\documentclass{amsart}
}

% For dealing with references we use the comment environment
\usepackage{verbatim}
\newenvironment{reference}{\comment}{\endcomment}
%\newenvironment{reference}{}{}
\newenvironment{slogan}{\comment}{\endcomment}
\newenvironment{history}{\comment}{\endcomment}

% For commutative diagrams we use Xy-pic
\usepackage[all]{xy}

% We use 2cell for 2-commutative diagrams.
\xyoption{2cell}
\UseAllTwocells

% We use multicol for the list of chapters between chapters
\usepackage{multicol}

% This is generally recommended for better output
\usepackage{lmodern}
\usepackage[T1]{fontenc}

% For cross-file-references
\usepackage{xr-hyper}

% Package for hypertext links:
\usepackage{hyperref}

% For any local file, say "hello.tex" you want to link to please
% use \externaldocument[hello-]{hello}
\externaldocument[introduction-]{introduction}
\externaldocument[conventions-]{conventions}
\externaldocument[sets-]{sets}
\externaldocument[categories-]{categories}
\externaldocument[topology-]{topology}
\externaldocument[sheaves-]{sheaves}
\externaldocument[sites-]{sites}
\externaldocument[stacks-]{stacks}
\externaldocument[fields-]{fields}
\externaldocument[algebra-]{algebra}
\externaldocument[brauer-]{brauer}
\externaldocument[homology-]{homology}
\externaldocument[derived-]{derived}
\externaldocument[simplicial-]{simplicial}
\externaldocument[more-algebra-]{more-algebra}
\externaldocument[smoothing-]{smoothing}
\externaldocument[modules-]{modules}
\externaldocument[sites-modules-]{sites-modules}
\externaldocument[injectives-]{injectives}
\externaldocument[cohomology-]{cohomology}
\externaldocument[sites-cohomology-]{sites-cohomology}
\externaldocument[dga-]{dga}
\externaldocument[dpa-]{dpa}
\externaldocument[sdga-]{sdga}
\externaldocument[hypercovering-]{hypercovering}
\externaldocument[schemes-]{schemes}
\externaldocument[constructions-]{constructions}
\externaldocument[properties-]{properties}
\externaldocument[morphisms-]{morphisms}
\externaldocument[coherent-]{coherent}
\externaldocument[divisors-]{divisors}
\externaldocument[limits-]{limits}
\externaldocument[varieties-]{varieties}
\externaldocument[topologies-]{topologies}
\externaldocument[descent-]{descent}
\externaldocument[perfect-]{perfect}
\externaldocument[more-morphisms-]{more-morphisms}
\externaldocument[flat-]{flat}
\externaldocument[groupoids-]{groupoids}
\externaldocument[more-groupoids-]{more-groupoids}
\externaldocument[etale-]{etale}
\externaldocument[chow-]{chow}
\externaldocument[intersection-]{intersection}
\externaldocument[pic-]{pic}
\externaldocument[weil-]{weil}
\externaldocument[adequate-]{adequate}
\externaldocument[dualizing-]{dualizing}
\externaldocument[duality-]{duality}
\externaldocument[discriminant-]{discriminant}
\externaldocument[derham-]{derham}
\externaldocument[local-cohomology-]{local-cohomology}
\externaldocument[algebraization-]{algebraization}
\externaldocument[curves-]{curves}
\externaldocument[resolve-]{resolve}
\externaldocument[models-]{models}
\externaldocument[functors-]{functors}
\externaldocument[equiv-]{equiv}
\externaldocument[pione-]{pione}
\externaldocument[etale-cohomology-]{etale-cohomology}
\externaldocument[proetale-]{proetale}
\externaldocument[relative-cycles-]{relative-cycles}
\externaldocument[more-etale-]{more-etale}
\externaldocument[trace-]{trace}
\externaldocument[crystalline-]{crystalline}
\externaldocument[spaces-]{spaces}
\externaldocument[spaces-properties-]{spaces-properties}
\externaldocument[spaces-morphisms-]{spaces-morphisms}
\externaldocument[decent-spaces-]{decent-spaces}
\externaldocument[spaces-cohomology-]{spaces-cohomology}
\externaldocument[spaces-limits-]{spaces-limits}
\externaldocument[spaces-divisors-]{spaces-divisors}
\externaldocument[spaces-over-fields-]{spaces-over-fields}
\externaldocument[spaces-topologies-]{spaces-topologies}
\externaldocument[spaces-descent-]{spaces-descent}
\externaldocument[spaces-perfect-]{spaces-perfect}
\externaldocument[spaces-more-morphisms-]{spaces-more-morphisms}
\externaldocument[spaces-flat-]{spaces-flat}
\externaldocument[spaces-groupoids-]{spaces-groupoids}
\externaldocument[spaces-more-groupoids-]{spaces-more-groupoids}
\externaldocument[bootstrap-]{bootstrap}
\externaldocument[spaces-pushouts-]{spaces-pushouts}
\externaldocument[spaces-chow-]{spaces-chow}
\externaldocument[groupoids-quotients-]{groupoids-quotients}
\externaldocument[spaces-more-cohomology-]{spaces-more-cohomology}
\externaldocument[spaces-simplicial-]{spaces-simplicial}
\externaldocument[spaces-duality-]{spaces-duality}
\externaldocument[formal-spaces-]{formal-spaces}
\externaldocument[restricted-]{restricted}
\externaldocument[spaces-resolve-]{spaces-resolve}
\externaldocument[formal-defos-]{formal-defos}
\externaldocument[defos-]{defos}
\externaldocument[cotangent-]{cotangent}
\externaldocument[examples-defos-]{examples-defos}
\externaldocument[algebraic-]{algebraic}
\externaldocument[examples-stacks-]{examples-stacks}
\externaldocument[stacks-sheaves-]{stacks-sheaves}
\externaldocument[criteria-]{criteria}
\externaldocument[artin-]{artin}
\externaldocument[quot-]{quot}
\externaldocument[stacks-properties-]{stacks-properties}
\externaldocument[stacks-morphisms-]{stacks-morphisms}
\externaldocument[stacks-limits-]{stacks-limits}
\externaldocument[stacks-cohomology-]{stacks-cohomology}
\externaldocument[stacks-perfect-]{stacks-perfect}
\externaldocument[stacks-introduction-]{stacks-introduction}
\externaldocument[stacks-more-morphisms-]{stacks-more-morphisms}
\externaldocument[stacks-geometry-]{stacks-geometry}
\externaldocument[moduli-]{moduli}
\externaldocument[moduli-curves-]{moduli-curves}
\externaldocument[examples-]{examples}
\externaldocument[exercises-]{exercises}
\externaldocument[guide-]{guide}
\externaldocument[desirables-]{desirables}
\externaldocument[coding-]{coding}
\externaldocument[obsolete-]{obsolete}
\externaldocument[fdl-]{fdl}
\externaldocument[index-]{index}

% Theorem environments.
%
\theoremstyle{plain}
\newtheorem{theorem}[subsection]{Theorem}
\newtheorem{proposition}[subsection]{Proposition}
\newtheorem{lemma}[subsection]{Lemma}

\theoremstyle{definition}
\newtheorem{definition}[subsection]{Definition}
\newtheorem{example}[subsection]{Example}
\newtheorem{exercise}[subsection]{Exercise}
\newtheorem{situation}[subsection]{Situation}

\theoremstyle{remark}
\newtheorem{remark}[subsection]{Remark}
\newtheorem{remarks}[subsection]{Remarks}

\numberwithin{equation}{subsection}

% Macros
%
\def\lim{\mathop{\mathrm{lim}}\nolimits}
\def\colim{\mathop{\mathrm{colim}}\nolimits}
\def\Spec{\mathop{\mathrm{Spec}}}
\def\Hom{\mathop{\mathrm{Hom}}\nolimits}
\def\Ext{\mathop{\mathrm{Ext}}\nolimits}
\def\SheafHom{\mathop{\mathcal{H}\!\mathit{om}}\nolimits}
\def\SheafExt{\mathop{\mathcal{E}\!\mathit{xt}}\nolimits}
\def\Sch{\mathit{Sch}}
\def\Mor{\mathop{\mathrm{Mor}}\nolimits}
\def\Ob{\mathop{\mathrm{Ob}}\nolimits}
\def\Sh{\mathop{\mathit{Sh}}\nolimits}
\def\NL{\mathop{N\!L}\nolimits}
\def\CH{\mathop{\mathrm{CH}}\nolimits}
\def\proetale{{pro\text{-}\acute{e}tale}}
\def\etale{{\acute{e}tale}}
\def\QCoh{\mathit{QCoh}}
\def\Ker{\mathop{\mathrm{Ker}}}
\def\Im{\mathop{\mathrm{Im}}}
\def\Coker{\mathop{\mathrm{Coker}}}
\def\Coim{\mathop{\mathrm{Coim}}}

% Boxtimes
%
\DeclareMathSymbol{\boxtimes}{\mathbin}{AMSa}{"02}

%
% Macros for moduli stacks/spaces
%
\def\QCohstack{\mathcal{QC}\!\mathit{oh}}
\def\Cohstack{\mathcal{C}\!\mathit{oh}}
\def\Spacesstack{\mathcal{S}\!\mathit{paces}}
\def\Quotfunctor{\mathrm{Quot}}
\def\Hilbfunctor{\mathrm{Hilb}}
\def\Curvesstack{\mathcal{C}\!\mathit{urves}}
\def\Polarizedstack{\mathcal{P}\!\mathit{olarized}}
\def\Complexesstack{\mathcal{C}\!\mathit{omplexes}}
% \Pic is the operator that assigns to X its picard group, usage \Pic(X)
% \Picardstack_{X/B} denotes the Picard stack of X over B
% \Picardfunctor_{X/B} denotes the Picard functor of X over B
\def\Pic{\mathop{\mathrm{Pic}}\nolimits}
\def\Picardstack{\mathcal{P}\!\mathit{ic}}
\def\Picardfunctor{\mathrm{Pic}}
\def\Deformationcategory{\mathcal{D}\!\mathit{ef}}
\usepackage{cleveref}
\crefname{originalsection}{original section}{original sections}
\Crefname{originalsection}{Original section}{Original sections}
\makeatletter
\newlabel{section-weakly-adic}{{0GXD}{source}{}{sourcepointer-1}{}}
\newlabel{section-weakly-adic@cref}{{[originalsection][0][]0GXD}{[1][0][]source}{}{}{}}
\makeatother
\title{Implications among the types of affine formal algebraic spaces}
\author{The Stacks Project Authors}
\date{stacks-project:formal-spaces.tex:local-snapshot-20261012; GNU FDL 1.2 or later}
\begin{document}
\maketitle
\noindent\textit{Editorial note.} All delivered mathematical text is the author's own, copied from the pinned Stacks Project chapter file formal-spaces.tex; no statement or proof step was rewritten, reordered or removed.\par\medskip
\noindent\textit{Editorial note.} The proof displays one description of a limit of rings; it lies inside the counted span and is reproduced verbatim.\par\medskip
\noindent\textit{Editorial note.} The proof contains the author's own wording 'there exits in particular'; it lies inside the counted span and is reproduced verbatim, as the delivered text is the author's.\par\medskip
\noindent\textit{Editorial note.} The counted spans carry 3 navigation stubs; no bibliography citation occurs in the spans.\par\medskip
\noindent\textit{Editorial note.} The preamble span points at the chapter's own preamble.tex shipped with the frozen source, used verbatim as staged.\par\medskip
\noindent\textit{Editorial note.} The delivered context span for the definition of a weakly adic ring carries a citation to [Gabber-Ramero]. Since the chapter ships no bibliography, that entry is attached to this card as bibliography.tex, copied character for character from the Stacks Project's own my.bib, so that the citation resolves without any entry being invented.\par\medskip
\noindent\textit{Allowed standard prerequisites (editorial).}\begin{itemize}
\item The types of affine formal algebraic spaces and of complete weakly admissible rings. Same chapter, the definitions of the six types appearing in the statement, used throughout. Fix the notions whose implications are asserted.
\item The implications between the corresponding types of topological rings. Same chapter, Section section-weakly-adic, invoked at proof lines 2469-2474. Give the first four parts once the space is written as a formal spectrum.
\item A weakly adic ring has a countable fundamental system of open ideals. Same chapter, Lemma lemma-weakly-adic at proof lines 2478-2481 and Definition definition-weakly-adic, invoked there. Gives part five.
\item Affine formal algebraic spaces as colimits of thickenings of affine schemes, with McQuillan meaning a surjection onto each stage. Same chapter, Definition definition-affine-formal-algebraic-space and the defining property of the McQuillan type, used at proof lines 2485-2490. Provide the presentation and the condition verified in part six.
\end{itemize}\medskip
\ifdefined\hypertarget\hypertarget{sourcepointer-1}{}\fi
\noindent\textit{Source reference.} Original section 0GXD: sources/stacks-formal-spaces/pinned-original-source.tex, source lines 1588--1588. Pointer outside the counted span, invoked once for context; the load-bearing argument of the counted statement is the author's proof reproduced in full in the proof span.\par\medskip
\begin{definition}
\label{definition-weakly-adic}
\begin{reference}
\cite[Definition 8.3.8]{Gabber-Ramero}
\end{reference}
Let $A$ be a linearly topologized ring.
\begin{enumerate}
\item We say $A$ is {\it weakly pre-adic}\footnote{In \cite{Gabber-Ramero} the
authors say $A$ is {\it $c$-adic}.} if there exists an ideal
$I \subset A$ such that the closure of $I^n$ is open for all $n \geq 0$
and these closures form a fundamental system of open ideals.
\item We say $A$ is {\it weakly adic} if $A$ is weakly pre-adic
and complete\footnote{By our conventions this includes separated.}.
\end{enumerate}
\end{definition}

\begin{lemma}
\label{lemma-implications-between-types}
Let $X$ be an affine formal algebraic space over a scheme $S$.
\begin{enumerate}
\item If $X$ is Noetherian, then $X$ is adic*.
\item If $X$ is adic*, then $X$ is adic.
\item If $X$ is adic, then $X$ is weakly adic.
\item If $X$ is weakly adic, then $X$ is classical.
\item If $X$ is weakly adic, then $X$ is countably indexed.
\item If $X$ is countably indexed, then $X$ is McQuillan.
\end{enumerate}
\end{lemma}

\begin{proof}
Statements (1), (2), (3), and (4) follow by writing $X = \text{Spf}(A)$
and where $A$ is a weakly admissible (hence complete) linearly
topologized ring and using the implications between the various
types of such rings discussed in Section \ref{section-weakly-adic}.

\medskip\noindent
Proof of (5). By definition there exists a weakly adic topological ring $A$
such that $X = \colim \Spec(A/I)$ where the colimit is over the ideals
of definition of $A$. As $A$ is weakly adic, there exits in particular
a countable fundamental system $I_\lambda$ of open ideals, see
Definition \ref{definition-weakly-adic}.
Then $X = \colim \Spec(A/I_n)$ by definition of $\text{Spf}(A)$.
Thus $X$ is countably indexed.

\medskip\noindent
Proof of (6). Write $X = \colim X_n$ for some system
$X_1 \to X_2 \to X_3 \to \ldots$ of thickenings of affine
schemes over $S$. Then
$$
A = \lim \Gamma(X_n, \mathcal{O}_{X_n})
$$
surjects onto each $\Gamma(X_n, \mathcal{O}_{X_n})$ because the transition
maps are surjections as the morphisms $X_n \to X_{n + 1}$ are closed
immersions. Hence $X$ is McQuillan.
\end{proof}

%% Bibliography retained so that the citation carried by the delivered context span of
%% S5-fsp-implications-between-types (lemma-implications-between-types) resolves.
%%
%% The delivered definition of a weakly adic ring cites [Gabber-Ramero].  The entry below is
%% copied character for character from the pinned Stacks Project bibliography file my.bib
%% (the @ARTICLE{Gabber-Ramero} record), which is the file the Stacks Project itself uses
%% to resolve this key.  Nothing was rewritten, reordered or invented, and the citation key
%% was not changed.
%%
%% Only this one key is cited by the card, and only this one entry is retained.
\begin{thebibliography}{99}

\bibitem[Gabber-Ramero]{Gabber-Ramero}
Gabber, Ofer and Ramero, Lorenzo.
\newblock Foundations for almost ring theory--Release 7.5.
\newblock arXiv preprint math/0409584, 2004.

\end{thebibliography}
\end{document}
